% =====================================================================
%  Harald Høffding, Benedikt Spinoza
%  In: Nær og Fjern.  Et Ugeblad, redigeret af P. Hansen.  Sjette Bind (Nr. 222-274).  Kjøbenhavn 1877: C. A. Reitzel; Otto B. Wroblewsky.  Nr. 233, Søndag den 17. December 1876, pp. [1]-3 (the weekly numbers each issue's pages from 1), the lead article.
%  Fraktur, two columns; French titles in antiqua (\textit{} here).
%
%  TRANSCRIPTION -- GENERATED; DO NOT EDIT BY HAND.  Høffding's four pieces of 1875-76 (Herbert Spencers Filosofi; the
%  review of Wilkens; Montaignes Betydning i Etikens Historie; Benedikt Spinoza) are built by ONE pipeline
%  (ebook-method/hoeffding-1875-76/) into a combined working file, which split.py cuts into four transcriptions.
%  Rebuild from the repo root:   sh ebook-method/hoeffding-1875-76/rebuild.sh
%
%  STATUS: COMPLETE (2026-10-15): the volume's title page and the whole text, pp. [1]-3.
%
%  EDITION.  Internet Archive naer-og-fjern-1876-77 (a private scan, 305 ppi colour JP2, uploaded 2024); extract ~/bibliotek/Høffding, Harald/1876-naer-og-fjern-6-nr233.pdf (SCANS.tsv, sha256).  No second copy found (Internet Archive, NB, Runeberg, HathiTrust).  The masthead (an engraved vignette) is described, not reproduced; the title and byline span both columns; the article is unsigned.
%
%  ACCURACY (2026-10-15; one collation for all four texts, 57 working pages, ~17,100 words).
%  (a) The DRAFT is our own tesseract reading (tessdata_best: dan for Det nittende Aarhundrede, script/Fraktur for Nær og
%      Fjern; 300 dpi) of the Archive's JP2 images, laid over the working PDF and turned into LaTeX by draft.py.  It was
%      chosen over the Archive's own OCR (tesseract 5.2/5.3, hOCR) by measurement against the Harvard copy's ABBYY OCR
%      (words unmatched: 3.0% against 9.5% in the Montaigne volume; the Archive reads aa as åa and long s as f).  The
%      Archive's OCR is the WITNESS: collated word by word (collate/check.py): 900 disagreements, 76 settled by check.py.
%  (b) Rules (rules.py: the Archive's å for a, long s read f, glyph faults against attested words, 326) and a THIRD reading
%      voted: the ABBYY OCR of ANOTHER COPY (Google's scan of the Harvard volume), by context and by word alignment
%      (rules.py harv-*, vote2.py hv-*), never on figures and never making a word the corpus does not attest.  AUDIT: 30
%      letter-changing rule/vote decisions read blind at the image: 0 contradicted (6 undetermined crops).
%  (c) The rest of Det nittende Aarhundrede read at the image BLIND (round b1: 210 crops on 12 sheets, five readers, 23/24
%      planted controls right; ~80 decided by hand from the readers' lines).  OPEN 0.
%  (d) Every FOOTNOTE read whole at the image (readers N1, N2: 26 notes on 23 pages), set from the reading; the note marks
%      checked at zooms (four the draft had misplaced, moved).  STRUCTURE read at the image (reader K): the title pages, the
%      heads, the endings and signatures, the verse on p. 116.
%  (e) Nær og Fjern (FRAKTUR, two columns, no second copy): each column PROOFREAD WHOLE at the image against the draft by
%      two readers (P1, P2; 3 of 4 planted errors found): 171 corrections (proofapply.py); then an unattested-word sweep
%      and a long-s / ø sweep, every doubtful word read at a zoom (proof-fixes.tsv, final-fixes.tsv).
%  (f) LETTERSPACING (\emph{} here) MEASURED at the image (spacemeasure.py, white runs between letters, threshold 5 px: 31
%      runs), every run judged at the image by reader K (13 kept in the body; the notes' spacing read by N1/N2).  Italics:
%      none found by stroke slant (check.py --italics).
%  (g) QUOTATION MARKS: the print quotes „...“ only (every quotation read at the image in b1, the notes and the structure);
%      both OCRs read „ as , or ,, and “ as ' ” " »: set by class (quotes.py), outside the notes.
%  (h) Sweeps on the collated text, every doubtful spot read at a zoom: unattested words (229 words in no other
%      transcription: 12 words mended, the misprints below confirmed); specks both OCRs had read as marks (specksweep.py,
%      marks the Harvard copy lacks mid-sentence: 16 removed).  The TRANSLATORS' odd spots (64) were read at zooms before
%      the translation was finished: 22 more faults mended (specks, ɔ: read as 2:, a running head and a footnote left in
%      the text, two paragraph breaks at page turns), 3 misprints confirmed (final-fixes.tsv).
%  (i) SAMPLED PAGE CHECK (2026-10-15, budget-limited to 3 pages): pp. 113 (Spencer), 217 (Montaigne) and one Nær og Fjern
%      column (p. 2, right), chosen at random (seed 20261015, one per stratum), transcribed BLIND at the image by one
%      reader (1,143 words) and compared word by word (pagecheck.py): 10 differences; decided at zooms: 3 reader slips;
%      7 errors of the collated text (2 words: „for“ for „før“, „skikkelig“ for the print's „stikkelig“; 5 marks: four
%      specks both OCRs had read as stops or quotes, one stop hidden by a planted control) = ~6 per 1000 words, all mended.
%      The specks led to specksweep.py over the whole of Det nittende Aarhundrede: 16 more found at zooms and removed.
%  PRINTED AS IS (misprints kept, marked at the line):
%    p. 2: „stilt“ for skilt (ſtilt); proofreader P1
%    p. 2: „paaden“ for paa den, set solid; sweep zoom
%    p. 2: „kartersianske“ for kartesianske; sweep zoom
%    p. 2: „stikkelig“ for skikkelig (ſtikkelig); page check, zoom
%    p. 3: „tehoretiske“ for theoretiske; proofreader P2, sweep zoom
%  CONVENTIONS.  Page turns: % --- p. N --- and \opage{N} (the printed page); Nær og Fjern's second column is marked % (column 2); a word divided over a page turn is
%  split with %.  Footnotes as \footnote at their marks, numbered per page as printed (*), **)).  \emph{} = letterspacing.
%  The running heads and the sheet signatures are not transcribed.  The ø rule (CLAUDE.md).
% =====================================================================
\documentclass[12pt,a4paper,openany]{book}
\usepackage[utf8]{inputenc}
\usepackage[T1]{fontenc}
\usepackage{libertinus}
\usepackage{libertinust1math}
\usepackage[danish]{babel}
\usepackage[protrusion=true,expansion=true]{microtype}
\usepackage{setspace}
\usepackage[margin=1.3in]{geometry}
\usepackage{fancyhdr}
\usepackage{hyperref}
\hypersetup{pdftitle={Benedikt Spinoza}, pdfauthor={Harald Høffding}, hidelinks}
\pagestyle{fancy}
\fancyhf{}
\fancyhead[LE,RO]{\thepage}
\fancyhead[C]{\textit{H. Høffding: Benedikt Spinoza (1876)}}
\fancypagestyle{plain}{\fancyhf{}\fancyhead[LE,RO]{\thepage}}
\setstretch{1.3}
\usepackage{marginnote}
\usepackage{graphicx}
\DeclareUnicodeCharacter{0254}{\reflectbox{c}}   % ɔ (id est)
\makeatletter
\renewcommand{\thefootnote}{\ifcase\value{footnote}\or *\or **\or ***\else\arabic{footnote}\fi)}   % the print marks each page's notes *), **)
\makeatother
\newcommand{\opage}[1]{\marginnote{\footnotesize\textit{[#1]}}}
\newcommand{\signature}[1]{\par\smallskip\hfill #1\hspace*{2em}\par}
\pdfstringdefDisableCommands{\renewcommand{\footnote}[2][]{}\renewcommand{\opage}[1]{}}

\begin{document}

\begin{titlepage}
\begin{center}
{\LARGE\bfseries Nær og Fjern.}\\[0.5em]
\rule{3em}{0.4pt}\\[0.5em]
{\large \emph{Et Ugeblad},}\\[0.3em]
{\small\bfseries redigeret af P. Hansen.}\\[0.5em]
\rule{3em}{0.4pt}\\[0.8em]
{\small \emph{Sjette Bind}.}\\[0.3em]
{\small\bfseries (Nr. 222---274.)}\\
\vfill
{\small\bfseries Kjøbenhavn 1877.}\\[0.4em]
{\small\bfseries C. A. Reitzel.\qquad Otto B. Wroblewsky.}
\end{center}
\end{titlepage}
\clearpage

% --- p. 1 ---
\setcounter{footnote}{0}%
\opage{1}\begin{center}{\small Kjøbenhavn.  Udkommer hver Søndag.\quad {\large\bfseries Nr. 233.}\quad Den 17. December 1876.  Abonnementspris: 3 Kr. 50 Øre Kvartalet.}\\[0.3em]\rule{\textwidth}{0.4pt}\\[0.8em]{\Large \emph{Benedikt Spinoza.}}\\[0.4em]{\small Af \textbf{Dr. phil.} Harald Høffding.}\\[0.4em]\rule{3em}{0.4pt}\end{center}\medskip

Vor Tids Iver for at oprejse Statuer og Mindesmærker
har ofte bragt mig til at tænke paa Ordet om dem, der
kalke Profeternes Grave. Man kan ikke holde den vemodige
Tanke borte, at Anerkjendelsen af Fortidens Storhed
ikke hindrer Miskjendelse af Samtidige. Men det er
paa den anden Side naturligt, at en friere og forsonligere
Betragtning lettest kommer frem overfor det Forbigangne,
hvor Alt er afsluttet, Kampen endt og de Kæmpendes
Stilling og Motiver belyste ved deres Plads og Vilkaar
indenfor den historiske Udvikling. For Spinozas Vedkommende
synes det i hvert Tilfælde at forholde sig saa.
Selv de, der ikke kunne sympathisere med hans Standpunkt
og hans Lære, kunne ikke nægte ham deres Beundring for
hans ædle og milde Karakter og maa indrømme hans store
Betydning for den senere Videnskab og Litteratur. Og
naar man erindrer, hvor megen Forfølgelse og Miskjendelse
Spinoza maatte lide, hvor ene han stod i sin Tid,
og hvor fremmed denne følte sig overfor hans Ideer, saa
at der hengik over 100 Aar, før de kunde trænge igjennem,
saa vil den Statue, som nu skal oprejses for ham i Haag
ved Bidrag ikke blot fra hans Fødeland, men ogsaa fra
det øvrige Europa, staa som et Vidnesbyrd om, at der er
foregaaet et Fremskridt ikke blot i Forstaaelse af det Store
og Betydningsfulde, men ogsaa i frisindet og tolerant
Tænkemaade.

At han stod ene i sin Tid, hænger maaske noget sammen
med, at han oprindelig var fremmed for den nyere
% (column 2)
\setcounter{footnote}{0}%
vesterlandske Videnskab og Litteratur. Hans Forældre vare
portugisiske Jøder, som havde forladt Fødelandet for deres
Troes Skyld og i Holland, Religionsfrihedens Vugge, havde
søgt sig et nyt Hjem. I Amsterdam var der paa den Tid
en blomstrende jødisk Menighed med en Synagoge og en
Rabbinskole, ved hvilken navnkundige Lærere virkede. I denne By fødtes \emph{Benedikt Spinoza} den 24. November
1632. Paa Grund af den Begavelse, han tidlig lagde for
Dagen, bestemtes han til den lærde Stand, det vil sige,
til at blive Rabbiner. Han gjorde ogsaa saa raske Fremskridt
paa denne Bane, at han allerede i sit 16. Aar var
en anseet Talmudist. Men snart begyndte hans skarpe
Forstand og utrættelige Forskerlyst at føle sig utilfreds
indenfor Traditionens og Theologiens Skranker. Rimeligvis
var det paa denne Tid, at han lærte Latin og begyndte
at studere de klassiske Forfattere, Naturvidenskaberne
og den kartesianske Filosofi. Hans Lærer i Humaniora
var Lægen van Ende, der havde en Slags Undervisningsanstalt
i sit Hus og var bekjendt for sine frie religieuse
Anskuelser. Et Sagn fortæller, at Spinoza blev forelsket
i hans Datter, der hjalp Faderen med Undervisningen,
men at hun foretrak en Meddiscipel for ham. Dog har
man i den nyeste Tid forkastet denne Beretning, da van
Endes Datter efter de Oplysninger, som ere fundne, paa
den Tid, hvorom her er Tale, kun var lidt over ti Aar
gammel.

De Tvivl, der rejste sig hos ham, og de nye Tanker,
% --- p. 2 ---
\setcounter{footnote}{0}%
\opage{2}som begyndte at forme sig, have sikkert tidlig fjernet ham
noget fra hans Troesfæller og gjort hans Besøg i Synagogen
mindre hyppige; thi der fortælles, at man lovede
ham en ikke lille aarlig Pension, naar han blot vilde besøge
Synagogen af og til. Ja, man skal endog have forsøgt
voldsommere Midler imod ham. En Aften, da han
forlod Synagogen, overfaldt En ham med en Dolk, hvis
Stød dog kun flængede hans Klædning. Man havde
Mistanke om, at han vilde træde over til Kristendommen,
og vilde nødig afstaa en saa begavet Tilhænger til en
anden Konfession. Dog maa det ogsaa fremhæves, at
Tilhængerne af den kartersianske Filosofi paaden Tid bleve % PRINTED AS IS (paaden: for paa den, set solid; sweep zoom) % PRINTED AS IS (kartersianske: for kartesianske; sweep zoom)
betragtede med Mistanke ogsaa indenfor den kristne Kirke.
Den anden Dortrechtersynode havde nylig erklæret dem for
udelukkede fra kirkelige Embeder. Den lille jødiske Menighed,
som maatte vaage endnu strengere over sine Medlemmer,
greb, da hverken gode eller onde Midler frugtede,
tilsidst den Udvej, at udstøde den Frafaldne fra sin Midte.
Dette skete i Juli Maaned 1656. Banformularen blev
for nogle Aar siden funden i den jødiske Menigheds Arkiv
i Amsterdam; den er holdt i det gamle Testamentes Stil
og benytter til Dels dets Ord. „Forbandet“ --- hedder
det bl. A. --- „være han om Dagen, forbandet være han
om Natten, forbandet, naar han lægger sig, forbandet,
naar han rejser sig fra sit Leje! Forbandet være hans
Udgang, forbandet hans Indgang!“ „Vider, at Ingen
maa have noget at skaffe med ham hverken mundtlig eller
skriftlig, Ingen maa vise ham nogen Gunst, Ingen dvæle
under hans Tag, Ingen nærme sig ham paa mere end fire
Alen, Ingen læse Noget, han har forfattet!“

I sit 24. Aar var Spinoza saaledes stødt ud i Verden,
voldsomt stilt fra Familie og Stammefrænder, og % PRINTED AS IS (stilt: for skilt (ſtilt); proofreader P1)
uden Tilhold noget andet Sted. Det kom ham nu til
Gode, at Talmud paalagde Enhver, der uddannede sig til
Rabbiner, at lære et Haandværk. Han havde lært den
Kunst at slibe optiske Glas og havde bragt det til stor
Færdighed heri. Derved skaffede han sig nu sit tarvelige
Udkomme. Men han fik endnu ikke Lov til at leve i Ro.
Paa den jødiske og den kristne Gejstligheds Opfordring
forviste Amsterdams Øvrighed ham i Aaret 1660 fra Staden
„for at hævde Ordenen og Subordinationen.“ Han
boede i de nærmest følgende Aar i Rhynsburg, en lille
Flække ved Leiden, hvor en arminiansk Sekt, de saa kaldte
Kollegianter, havde fundet et Fristed. Senere flyttede han
til Voorburg ved Haag, og endelig (1669) paa sine Venners
Opfordring til selve Haag.

Om Spinozas Karakter og indre Liv i hans Ungdom
veed man kun, hvad man kan slutte af de ydre Begivenheder
og af Brevene og Skrifterne fra hans senere Aar.
Hans Biografer slutte gjerne tilbage fra den Ro og Resignation,
som udpræger sig i disse, og mene, at han ogsaa
i sin Ungdom har lagt saadanne Egenskaber for Dagen;
dog mene Nogle, at man kun kan tænke sig dem hos den
modne Tænker som Resultat af lidenskabelige Kampe i
Ungdommen. Det er i hvert Tilfælde klart, at han tidlig
har sammenfattet al sin Higen og Tragten i een Attraa ---
den at erhverve en klar og gyldig Erkjendelse af Menneskets
Væsen og Forhold til hele Naturen. Hele hans
Karakter har et intellektuelt Præg; Undersøgelse og Betragtning
er for ham det Højeste. Han skildrer selv i sin
Afhandling om „Forstandens Forbedring“, hvorledes de
% (column 2)
\setcounter{footnote}{0}%
Goder, de Fleste søge, Ære, Rigdom og sandselig Nydelse,
havde tabt deres Værd for ham, og det nu gik op for
ham, at det sande og højeste Gode var „Erkjendelsen af
vor Aands Sammenhæng med hele Naturen.“ Han søgte
i Sandhedserkjendelsen Mere end blot Theori; han søgte
en Livsanskuelse. Og hans Storhed bestaaer her netop i
den Harmoni, der hos ham var mellem Erkjendelsen og
Karakteren, Tanken og Livet.

Efter Udgivelsen af hans første Skrift, en Fremstilling
af den kartesianske Filosofi med Antydninger af de Punkter,
hvor den maatte føres videre, blev han snart berømt.
Mange opsøgte ham i hans Bolig, og han førte en betydelig
Brevvexling. En af hans Korrespondenter, Oldenburg,
den første Præsident for Videnskabernes Selskab i
London; nævner som det Tiltrækkende hos ham „grundig
Viden i Forbindelse med Mildhed og fin Dannelse.“ I Aaret
1670 udgav han en „theologisk-politisk Afhandling“, som gik
ud paa at vise, „at Tænkefriheden ikke blot kan indrømmes
uden Skade for Fromheden og for Statens Fred, men at
den end ikke kan ophæves, uden at disse med det Samme
gaa til Grunde“. Han forlangte her Ret for Enhver til
at fortolke de religieuse Lærdomme efter bedste Indsigt ---
og ikke blot danne sig sin egen Overbevisning, men ogsaa
udtale den. Han anfører (idet han synes at glemme,
hvorledes det var gaaet ham selv) sin Fødeby, Amsterdam,
som Exempel paa, at en saadan Frihed kan bestaa med
et ordnet Samfunds Væsen. „I denne blomstrende og
udmærkede Stad leve Mennesker af alle Nationer og
Troesbekjendelser i største Endrægtighed. Der er ingen nok
saa forhadt Sekt, uden at dens Tilhængere jo (naar de
blot Ingen krænke, give Enhver Sit og leve stikkelig) beskyttes % PRINTED AS IS (stikkelig: for skikkelig (ſtikkelig); page check, zoom)
af den offentlige Statsmagt“. Han er den Første,
der har hævdet Tolerancens Princip i denne fulde Udstrækning.
Det anførte Værk indeholdt desuden historisk-kritiske
Undersøgelser af det gamle Testamentes Bøger,
hvorved Grunden blev lagt til hele den nyere Bibelkritik,
og rationalistiske Fortolkninger af de vigtigste Dogmer.
Den Forargelse, dette Skrift vakte, afskrækkede Spinoza
fra at udgive Mere. Hans Hovedværk, „Ethiken“, cirkulerede
i Afskrifter mellem hans Venner og udkom først
efter hans Død.

I hele den sidste Del af sit Liv led han af en
tærende Brystsygdom. Han fortsatte dog stadig sine Studier,
og med saadan Iver, at han en Gang i tre Maaneder
ikke forlod sit Hus. Han boede hos Maleren van Spyck,
i hvis Familie han ofte tilbragte sine Hviletider; han talte
navnlig gjerne med sin Vært og hans Kone om Religionen
og om de Prædikener, de havde hørt i Kirken. Endnu Dagen før sin Død var han nede hos dem for at ryge en
Pibe Tobak og tale om Dagens Prædiken. Han formanede
stedse Folk til at blive i deres Konfession og fremhævede
det praktisk-ethiske Moment i Religionen.

Spinoza døde d. 21. Februar 1677. Ingen var tilstede
hos ham uden hans Ven Lægen Ludvig Meyer.
Der udbredte sig Rygter om, at han skulde være død i
Fortvivlelse eller endog have begaaet Selvmord, ligesom
der ogsaa om hans Karakter og Adfærd vare meget slemme
Rygter i Omløb. En luthersk Gejstlig, Colerus, der kom
til at bo i samme Hus, som Spinoza tidligere havde boet
i, tog sig for at undersøge Sammenhængen og udgav en
Levnetsbeskrivelse af Spinoza (paa Hollandsk, oversat paa
% --- p. 3 ---
\setcounter{footnote}{0}%
\opage{3}Fransk: \textit{«La vie de Spinoza, tirée des écrits de ce fameux philosophe et du témoignage de plusieurs personnes dignes de foi, qui l’ont connu particulièrement, par Jean Colerus. A la Haye 1706»} --- det bedste Kildeskrift
til Spinozas Levnet) og han kunde modsige alle
de Rygter, som Fordommen havde avlet og Fanatismen
udspredt. Uagtet Colerus selv fordømmer Spinozas Lærdomme,
fremhæver han hans rene Sæder, milde Karakter
og store Uegennyttighed. Denne sidste Egenskab udviste
Spinoza navnlig overfor sine Søstre, hvem han overlod
sin Arvepart, ligesom han flere Gange afviste Tilbud om
Pengeunderstøttelse fra sine Venner. Et Professorat i
Heidelberg, som blev ham tilbudt, afslog han for ej at
hindres i sine rolige Studier og af Frygt for Intolerancen.

Man har ofte vel udelukkende udledt Spinozas Lære
af Kartesianismen. Af et for nogle Aar tilbage fundet
Ungdomsskrift fremgaaer det, at han har været stærkt paavirket
af Middelalderens jødiske Religionsfilosofer og af
Giordano Bruno, Renaissancens største Filosof, der hævdede
Naturens Uendelighed og Enhed. Hvad Spinoza
fandt hos Descartes, var den klare, mathematiske Methode
og en mekanisk Naturopfattelse. Udgaaende fra Ideen om
den uendelige Substans, der paa een Gang er Natur og
Gud, søger han at vise, at det er den, der rører sig i
Alt. Den materielle og den aandelige Verden ere forskjellige
Former for Aabenbaringen af den uendelige
Enhed.

Slige Tanker maatte i Følge Sagens Natur længe
ligge hen uforstaaede; baade religieuse og filosofiske Fordomme hindrede Tilegnelsen af dem. Spinoza laa, efter
Lessings Ord, som en død Hund, haanet og sparket af
alle Parter. Da kom hans Opstandelse i Slutningen af
det 18. Aarhundrede med den friere og friskere Bevægelse
i Poesi, Filosofi og Videnskab. Det vilde føre for vidt
her at udvikle, hvad han har været for Aander som Lessing,
Herder, Goethe, Novalis, Schelling og Schleiermacher.
For Goethe aandede der i Spinozas Skrifter en „Fredens
Luft“; han lærte af dem at se, at der er en stor, nødvendig
Sammenhæng i Alt, og beundrede den dybe Resignation,
som udtaler sig i dem. I det Skrift af Schleiermacher, som indleder den religieuse Fornyelse efter det 18.
Aarhundredes Rationalisme, hylder han i en begejstret
Lovprisning Spinozas dybe og i Sandhed religieuse Aand.
I de Forelæsninger, som Schellings Discipel Henrik
Steffens ved dette Aarhundredes Begyndelse holdt i Kjøbenhavn,
og som havde saa mægtig en Indflydelse paa
vor Litteraturs Udvikling, finder man ligeledes en smuk
og begejstret Lovtale over Spinoza. „Aldrig kan jeg“, siger
Steffens bl. A., „tænke mig denne herlige Mand i sin
rolige Ensomhed, med et Lys i sin Sjæl, som intet Tilfælde
formaaede at bortjage, med en indvortes Fred, som
ingen Forfølgelser kunde forstyrre, uden at en Straale af
% (column 2)
\setcounter{footnote}{0}%
hans Lys vækker det skjønneste Haab hos mig.“ (Man
efterlæse hele Begyndelsen af den 2. Forelæsning).

Hvad man i Aarhundredets Begyndelse forsøgte ad
poetisk og spekulativ Vej --- at opfatte Naturen og Tilværelsen
i sin Helhed og Enhed, det arbejder i vore Dage
den reale Forsken henimod. I vor Tids Naturvidenskab
og Filosofi er der en stærk enhedssøgende Tendens. Der
aabner sig vide Horizonter, og der knyttes Forbindelser og
dannes Overgange, hvor man tidligere standsede ved et
gabende Svælg. Dette maa naturlig levendegjøre Mindet
om den ensomme hollandske Tænker i det 17. Aarhundrede.
Spinoza er ogsaa derved beslægtet med den moderne
Realisme, at han trods sin spekulative og deduktive Methode
hævder den strenge kausale Sammenhæng og fordrer
en mekanisk Forklaring gjennemført overalt. Særlig hvad
Forholdet mellem Sjæl og Legeme angaaer, synes hans
Theori at være den, mod hvilken de nyeste psykologiske og
fysiologiske Undersøgelser pege hen.

Vi kunne ikke længer bruge Spinozas Methode i
Filosofien. Hvad der for hans Tanke stod som det Første,
det sikre Udgangspunkt: det absolute Væsen, er for os
netop det, der evig unddrager sig Tanken. Vi maa søge
ad Erfaringens Vej --- som ogsaa for Spinoza var den
eneste, ad hvilken vi kunne lære Enkeltvæsenerne at kjende
--- at finde Lovene og Formerne for Aandens Liv og dens
Sammenhæng med hele Tilværelsen. Men Spinozismen
staaer som et Mønster ved sin Dristighed og Konsekvens;
den er et Forsøg, som i denne Form kun har sit Sidestykke
i Hegels System, og dette kan i Originalitet og i
Liniernes Renhed ikke maale sig med Spinozas. Og foruden
det Befrugtende og Dybsindige i Grundtanken er
der ogsaa hos Spinoza, som hos enhver stor Tænker,
underordnede Punkter af megen Interesse og blivende
Betydning. Jeg har allerede omtalt hans historiske
Kritik og hans Hævden af Tolerancens Princip, og skal
endnu kun nævne hans psykologiske Fremstilling af Lidenskaberne,
deres Oprindelse, Udvikling og Vexelspil. Det
er en i sin Art saa klassisk Fremstilling, at Johannes
Müller endnu for nogle og tredive Aar siden aftrykte den
hel og holden i sin berømte „Handbuch der Physiologie“.
Den har først i vor Tid fundet Sidestykker i de engelske
Psykologers Fremstillinger.

Men aldeles afseet fra de tehoretiske Resultater staaer % PRINTED AS IS (tehoretiske: for theoretiske; proofreader P2, sweep zoom)
Spinoza som et stort Forbillede paa den sande Forskers
Karakter, og ved den Mildhed og Ro, han bevarede
i sit Indre saa vel som overfor Andre, skjøndt han
havde lidt under Livsanskuelsens Brydninger, staaer han
særlig som en Skikkelse, af hvis Betragtning man, som Steffens sagde, kan hente „det skjønneste Haab“ i en Tid som vor, hvor Anskuelserne ikke brydes mindre, selv om
Brydningen maaske foregaaer under mindre voldsomme
Former.

\end{document}
